\chapter{تصحيح أخطاء الآلة}
\chaptermark{تصحيح أخطاء الآلة}
\label{sec:debugging}
\begin{chaptergoals}
\item لماذا يكفي مسبار يسمع ألف عصبون من أصل $8.6\cdot10^{10}$ لفكّ شفرة الحركة؛
\item لماذا يجب أن يستخرج الجهاز المزروع النبضات على الرقاقة ولا يرسل إلا الأحداث، لا التدفق الخام؛
\item كيف يقرأ مفكِّك الشفرة ترددات المجموعات، ولماذا لا تستخرج الواجهات إلا بضعة بتات في الثانية؛
\item لماذا الكتابة في الدماغ أصعب من القراءة، وما الذي لا يراه المسبار.
\end{chaptergoals}

\section{كيف تُصحَّح أخطاء آلة بلا مخطط}

المهندس الذي يُعطى لوحة تعمل بلا مخطط يصحح أخطاءها بأربع حيل. يضع \emph{مسبارًا} وينظر ما الذي يسري في السلك. \emph{ويحقن إشارة اختبار} وينظر ما الذي تغيّر. \emph{ويفصل} جزءًا من الدارة وينظر ما الذي توقف عن العمل. \emph{ويقرأ سجلات التحكم} ليعرف في أي وضع تعمل اللوحة. هذا الكتاب كله محاولة لقراءة مخطط الدماغ، وفي هذا الفصل سننظر في أيّ الحيل الأربع يُحسن المهندسون استعمالها اليوم، وما الذي تقوله لهم الفصول السابقة عمّا ينبغي أن يتوقعوه.

أبرز أمثلة هذا التصحيح اليوم واجهات الدماغ والحاسوب: أجهزة تقرأ نبضات العصبونات وتحوّلها إلى أوامر لآلات خارجية، أو تحقن في الدماغ إشارات من الخارج بالعكس. ولها يُخصَّص معظم الفصل، وأولًا لما تفعله شركة \foreignlanguage{english}{Neuralink}.

\section{المسبار: ماذا يسمع القطب}

مسبار الدماغ قطب رفيع مغروس في النسيج. يسمع نبضات العصبونات الواقعة في نصف قطر من رتبة مئة ميكرومتر، وهي عادةً من عصبون واحد إلى بضعة عصبونات. النبضة على القطب ومضة من عشرات إلى مئات الميكروفولتات تدوم نحو جزء من الألف من الثانية، ومن شكلها يمكن تمييز العصبونات المتجاورة بعضها من بعض. وأوضح ما يُسمع عصبونات \nt{GLUT} الكبيرة: فهي الأكثرية في القشرة (الجدول~\ref{tab:types})، وتياراتها أقوى.

والصعوبة الرئيسية ظاهرة على الفور. المسبار الجيد اليوم مئات أو آلاف الأقطاب، أما عدد العصبونات في الدماغ فـ$8.6\cdot10^{10}$. نحن نتنصّت على نحو جزء من مئة مليون من الآلة. فلماذا يمكن أن نفهم شيئًا أصلًا من عيّنة ضئيلة كهذه؟ أجاب الفصل~\ref{sec:code}. العصبونات المتجاورة تتشارك ضجيجها، فيصطدم المتوسط على المجموعة بحد (القضية~\ref{prop:pooling}): مئة عصبون تحمل تقريبًا ما تحمله ألف. ثم إن المسألة التي تحلها القشرة الحركية منخفضة الأبعاد: للذراع مفاصل قليلة (الفصل~\ref{sec:muscles})، ونشاط مئات العصبونات التي تتحكم فيها ينحصر في فضاء من عشرة أو عشرين متغيرًا مشتركًا فقط~\cite{Gallego2017}. والمسبار الذي يسمع مئة عصبون يسمع تقريبًا كل هذه المتغيرات. ولو كان الدماغ يخزن في كل عصبون رسالة مستقلة، لما أجدى هذا التنصّت شيئًا.

\section{تدفق البيانات: الضغط على المسبار نفسه}

يُخرج المسبار تدفقًا هائلًا. لرؤية شكل النبضة يجب رقمنة كل قطب نحو $20$ ألف مرة في الثانية بدقة نحو عشرة بتات. وألف قطب تعطي
\begin{empheq}[box=\keybox]{equation}
\begin{aligned}
R_{\text{raw}} \;&=\; N\,f_s\,b \;\approx\; 10^3\cdot2\cdot10^4\cdot10 \;\approx\; 2\cdot10^8\ \text{bit/s},\\
R_{\text{spk}} \;&\approx\; N\,r\,\log_2\frac{e}{r\,\Delta t} \;\approx\; 8\cdot10^4\ \text{bit/s}.
\end{aligned}
\label{eq:probedata}
\end{empheq}
التقدير الثاني هو سعة تدفق النبضات~\eqref{eq:spikeentropy} عند $r=10$ نبضات في الثانية ودقة $\Delta t=1\ms$: كل ما فيه فعلًا يشغل حيزًا أصغر بألفين وخمسمئة مرة. ولجهاز مزروع هذا فرق حاسم: لا يمكن إرسال التدفق الخام لاسلكيًا عبر الجلد بالطاقة التي يُسمح بها داخل الرأس دون تسخين النسيج. لذلك يجب على الجهاز المزروع أن يستخرج النبضات على الرقاقة نفسها بجوار الأقطاب، وألا يرسل إلى الخارج إلا الأحداث: رقم القناة ولحظة النبضة. في الوصف الأول لمنظومة \foreignlanguage{english}{Neuralink} لم يكن الأمر قد بلغ ذلك بعد: الرقاقات تضخّم الإشارات وترقمنها، والتدفق الخام كله يخرج عبر كابل، وبرنامج في الخارج يستخرج النبضات؛ أما الضغط على الرقاقة والاتصال اللاسلكي فقد ذُكرا هناك خطةً~\cite{Musk2019}.

هذه حيلة مألوفة. العين تضغط الإشارة مئة مرة قبل أن ترسلها عبر كابل طويل (الفصل~\ref{sec:sensors})؛ وكاميرا الأحداث لا ترسل إطارات بل تغيرات. والمسبار، كي يتسع له السلك، مضطر إلى أن يفعل ما يفعله الدماغ نفسه: أن يتكلم بالنبضات وألا يتكلم إلا عن الجديد.

\section{ماذا تفعل \foreignlanguage{english}{Neuralink}}

جهاز \foreignlanguage{english}{Neuralink} مسبار مصمم لهذه القيود~\cite{Musk2019}. بدل مصفوفة صلبة من الإبر، خيوط مرنة كثيرة أرفع من الشعرة، على طول كل منها أقطاب: في الوصف الأول للمنظومة حتى $3072$ قطبًا على $96$ خيطًا، $32$ على كل خيط، وفي الجهاز المزروع في البشر، بحسب ما أعلنته الشركة، نحو ألف قناة على $64$ خيطًا. الخيوط المرنة تؤذي النسيج أقل، وتتحمل أفضل أن الدماغ داخل الجمجمة يتحرك قليلًا طوال الوقت. يغرس الخيوطَ روبوتٌ يتجنب الأوعية. وفي الوصف الأول، كما سبق، كان التدفق الخام~\eqref{eq:probedata} يخرج عبر كابل. أما الجهاز المخصص للبشر، بحسب ما أعلنته الشركة، فمبني بطريقة أخرى: الغلاف مغطى بالجلد كليًا، والنبضات تُستخرج في الداخل، والبيانات تخرج لاسلكيًا، والطاقة تصل بلا أسلاك.

مهمة الجهاز الأول هي القراءة. توضع الخيوط في جزء القشرة الذي يخطط حركات الذراع، ويحوّل مفكِّكُ الشفرة النبضاتِ إلى حركة مؤشر على الشاشة. فيتحكم إنسان مشلول اليدين في الحاسوب بتخيّل الحركة. الفكرة نفسها ليست جديدة: الواجهات القائمة على مصفوفات صلبة من مئة قطب سمحت منذ زمن بالتحكم في المؤشر~\cite{Hochberg2006}، وبكتابة النص بتخيّل حركات اليد في الكتابة~\cite{Willett2021}، بل بتحويل محاولات الكلام إلى نص على الشاشة بسرعة نحو ستين كلمة في الدقيقة~\cite{Willett2023}. الجديد عند \foreignlanguage{english}{Neuralink} هو الهندسة: عدد القنوات، واللاسلكية، والغلاف المغلق، والغرس الروبوتي، أي محاولة صنع مسبار يمكن حمله سنوات خارج المختبر. وبحسب علمنا، نُشرت نتائج التجارب على البشر في الغالب من الشركة نفسها لا في مقالات محكَّمة؛ وقد أعلنت الشركة، من بين أمور أخرى، أن بعض الخيوط انزاح بعد الغرس فنقص عدد القنوات العاملة. وهذا إزعاج عادي في تصحيح الأخطاء: المسبار الذي يتحرك بالنسبة إلى اللوحة يسمع أسلاكًا أخرى.

الاتجاه الثاني للشركة هو الكتابة: جهاز للبصر يحقن إشارات في القشرة البصرية لإنسان فقد بصره. وعنه نتحدث أدناه، في قسم الكتابة.

\section{القراءة: مفكِّك الشفرة متلقّيًا}

مفكِّك شفرة الواجهة متلقٍّ بمعنى القضية~\ref{prop:reader}: فهو يقرر بنفسه أي جزء من الرسالة يقرأ. عمليًا، تقرأ كل الواجهات تقريبًا \emph{ترددات المجموعات}: تعدّ نبضات كل قناة في نوافذ من بضع عشرات من أجزاء الألف من الثانية، وتجمعها بأوزان مختارة بحيث تنتج الحركة المقصودة. هذا ترميز التردد للمجموعة من القسم~\ref{sec:rate}، ولم يُختر مصادفة: مع بضع نبضات في النافذة لا يتحدد تردد عصبون واحد، أما المجموع على مئة فيعمل.

كم من المعلومات يمكن استخراجه؟ الكتابة ”باليد المتخيَّلة“ سارت بسرعة نحو تسعين حرفًا في الدقيقة~\cite{Willett2021}، أي حرف ونصف في الثانية: وحتى بخمسة بتات للحرف يكون ذلك أقل من ثمانية بتات في الثانية. والتحكم في المؤشر، بحسب ما أعلنته \foreignlanguage{english}{Neuralink}، يعطي نحو عشرة بتات في الثانية. أما الحد الأعلى~\eqref{eq:spikeentropy} لعصبون \emph{واحد} فعشرات البتات في الثانية، ولمئة عصبون آلاف. الواجهة تستخرج من العصبونات المتنصَّت عليها جزءًا صغيرًا مما يمكن أن تحمله. وعنق الزجاجة ليس السلك، بل عدد المتغيرات المستقلة التي تحملها المجموعة (القضية~\ref{prop:pooling})، ومدى فهم المفكِّك للشفرة التي، كما رأينا في الفصل~\ref{sec:code}، لم تُفكّ كاملة بعد.

وثمة خاصية ثانية مألوفة من الفصل~\ref{sec:motorlearning}. المفكِّك والدماغ يتعلم كل منهما من الآخر. يرى الدماغ إلى أين ذهب المؤشر، فيتلقى خطأً ويعدّل أوامره، وهو التعلم الحركي نفسه عبر سلك الخطأ. والمهندسون يعيدون من حين إلى آخر اختيار أوزان المفكِّك، لأن العصبونات تحت الخيوط تتغير والوصلات في الشبكة يُعاد تعلّمها. فينتج متعلمان يتكيف كل منهما مع الآخر؛ ولكي لا يفلت هذا الزوج من الضبط، يحسن الفصل بين سرعتي تعلمهما: ليتعلم أحدهما بسرعة والآخر ببطء.

\begin{keyblock}
\begin{proposition}[لماذا يعمل مسبار على ألف عصبون]
\label{prop:probe}
الواجهة التي تتنصت على جزء ضئيل من العصبونات تستطيع التحكم في الحركة لأن نشاط المجموعة منخفض الأبعاد وضجيجها متشارك: بضع مئات من العصبونات تحمل تقريبًا كل المتغيرات المشتركة. وللسبب نفسه لا تستخرج الواجهة إلا بضعة بتات في الثانية، بقدر المتغيرات المستقلة في المسألة، لا بقدر ما يتسع له تدفق النبضات. عمل الواجهات مُقاس؛ أما مقدار تحسّن المفكِّك حين تُفهم الشفرة فسؤال مفتوح.
\end{proposition}
\end{keyblock}

\section{الكتابة: أصعب من القراءة}

حقن إشارة في الآلة أصعب من القراءة منها. القطب الذي يمرّ عبره تيار لا يثير عصبونًا واحدًا، بل كل العصبونات والأسلاك حوله، دون تمييز للنوع أو الإشارة. ولا يمكن به كتابة شفرة يهم فيها \emph{أيّ} عصبون أطلق و\emph{متى} (الفصل~\ref{sec:code}): لا يمكن إلا تحديد \emph{أين} و\emph{بأي شدة} تحديدًا خشنًا.

وهذا يكفي للبصر جزئيًا. يرى الإنسان التيار عبر قطب في القشرة البصرية نقطةً مضيئة في موضع محدد من مجال الرؤية؛ وحين أُضيئت النقاط معًا، حتى ست عشرة نقطة في آن، تعرّف إنسان فقد بصره على بعض الحروف وميّز حدود الأشياء~\cite{Fernandez2021}. هذه شفرة بالموضع، ويمكن كتابتها. لكن تذكّر الفصل~\ref{sec:sensors}: العين لا ترسل إلى الدماغ بكسلات، بل وصفًا مضغوطًا: الحواف، والتغيرات، والإضاءة والإعتام، على أسلاك منفصلة. كتابة ”نقاط الضوء“ في القشرة كتابةُ بكسلات في آلة تنتظر عند مدخلها شفرة مختلفة تمامًا. لذلك ما زال البصر الاصطناعي أخشن مما يمكن توقعه من عدد الأقطاب. وثمة إشارات اختبار لا تحتاج إلى قطب: الخداع البصري يحقن عبر العين مدخلًا غير مألوف ويُخرج الآلة من وضعها المعتاد، وكل خطأ يدل على آليته (القسم~\ref{sec:illusions}).

\section{ما لا يراه المسبار}

القطب يسمع ناقل العناوين: نبضات عصبونات \nt{GLUT}، وبدرجة أقل عصبونات \nt{GABA} الصغيرة. لكننا رأينا في الفصول~\babelsublr{\ref{sec:serocomputer}--\ref{sec:acet}} أن وضع عمل الآلة كلها تحدده سجلات البث: تراكيز \nt{SERO} و\nt{DOPA} و\nt{NORA} و\nt{ACET}. وهذه لا يسمعها القطب: العصبونات المصدر قليلة إلى حد الضآلة، وإشارتها ليست نبضة بجوار المسبار بل تركيز في حجم. ومصحح الأخطاء الذي يرى ناقل البيانات ولا يرى سجلات التحكم سيبقى متعجبًا دائمًا: المدخل الواحد سيعطي استجابات مختلفة، لأن $g$ أو $a$ أو $\tau_\gamma$ تغيّر في مكان ما. يمكن قياس التراكيز بمستشعرات كيميائية داخل النسيج مباشرة~\cite{Bunin1998}، لكن هذه المستشعرات لا تُستعمل بعد في الواجهات. ويبقى خارج مجال رؤية المسبار تمامًا المخرجُ الثاني البطيء للآلة: الهرمونات في الدم (الفصل~\ref{sec:chemoutput})، وهي تقاس في عينات الدم لا بالقطب.

ولا يرى المسبار الأوزان، وفيها بالذات تُخزَّن الذاكرة كلها تقريبًا وكل ما تعلّمته الآلة. لا يمكن إلا تخمينها من الطريقة التي تتغير بها الاستجابات. ولا يرى الألم كما يشعر به الإنسان: رأينا في الفصل~\ref{sec:pain} أن شدة إشارة الإنذار تحددها الآلة نفسها، ببوابات في الحبل الشوكي ومنظّمات من الأعلى، فلا يمكن أن يُقرأ من المستشعر مقدار الألم.

\section{الخلاصة: التصحيح والأسئلة المفتوحة}

\begin{table}[t]
\centering
\footnotesize
\setlength{\tabcolsep}{4pt}
\renewcommand{\arraystretch}{1.15}
\begin{tabular}{>{\raggedright}p{2.2cm}>{\raggedright}p{3.6cm}>{\raggedright}p{3.4cm}>{\raggedright\arraybackslash}p{4.0cm}}
\toprule
حيلة التصحيح & ماذا تفعل الواجهة & ماذا يفسر الكتاب & ما المعروف \\
\midrule
المسبار & يسمع مئات إلى آلاف العصبونات & انخفاض الأبعاد والضجيج المتشارك (الفصل~\ref{sec:code}) & مُقاس \\
الضغط على المسبار & يرسل أحداثًا بدل الإشارة الخام & سعة تدفق النبضات~\eqref{eq:spikeentropy} & محلول هندسيًا \\
القراءة & من ترددات المجموعات إلى حركة ونص وكلام & المفكِّك متلقٍّ، ترميز التردد للمجموعة & يعمل؛ بضعة بتات في الثانية \\
التعلم المشترك & الدماغ والمفكِّك يتكيف كل منهما مع الآخر & سلك الخطأ (الفصل~\ref{sec:motorlearning}) & مُلاحَظ؛ قواعد التكيف موضوع بحث \\
الكتابة & التيار يضيء نقاطًا في مجال الرؤية & ترميز الموضع يمكن كتابته، وترميز التوقيت لا (الفصل~\ref{sec:sensors}) & النقاط مُقاسة؛ البصر المفيد ليس بعد \\
قراءة السجلات & لا تكاد تُجرى & قنوات البث (الفصول~\babelsublr{\ref{sec:serocomputer}--\ref{sec:acet}}) & المستشعرات الكيميائية موجودة في التجارب \\
\bottomrule
\end{tabular}
\caption{تصحيح أخطاء الآلة: ما تستطيعه واجهات الدماغ والحاسوب وما تقوله فصول الكتاب عنها.}
\label{tab:debugging}
\end{table}

يلخص الجدول~\ref{tab:debugging} الفصل. تستطيع واجهات الدماغ والحاسوب اليوم أن تضع مسبارًا على ناقل العناوين وأن تقرأ منه بضعة بتات في الثانية، وكون ذلك يعمل أصلًا يفسره انخفاض الأبعاد والضجيج المتشارك من الفصل~\ref{sec:code}. أما الكتابة في الآلة فلا تستطيعها إلا خشنةً، لأن شفرة مدخلها مضغوطة وموجهة إلى أسلاك منفصلة. وسجلات التحكم والأوزان، وهي ما يجعل الآلة قابلة للتعلم والضبط، لا يكاد المسبار يراها بعد.

كل سؤال مفتوح في الفصل~\ref{sec:synthesis} هو أيضًا مسألة لمصحح الأخطاء. أيّ شفرة تُقرأ، سيتضح حين تصير المسابير أكثف ويتعلم المفكِّكون استعمال توقيت النبضات لا تردداتها فقط. وكيف تتعلم شبكة متعددة الطبقات، يمكن فحصه بوضع مسابير في عدة طبقات معًا ومراقبة تغير استجاباتها أثناء التعلم. وما الذي تنقله قنوات البث، سيتضح حين تضاف المسابير الكيميائية إلى الكهربائية. هذا الكتاب محاولة لقراءة مخطط الآلة من سلوكها ومن القوانين التي يعرفها كل مهندس. ومصححات الأخطاء التي تُبنى اليوم ستتيح فحص هذا المخطط بالمسبار.

\section{تمارين}

\begin{exercise}[\lvlA]
\label{ex:17:1}
\emph{ميزانية القناة اللاسلكية.} الإرسال عبر الجلد يكلف $1\,\text{nJ}$ للبت، ولا يجوز للمرسل أن ينفق أكثر من $10\,\text{mW}$. ما القدرة اللازمة للتدفق الخام~\eqref{eq:probedata} عند $N=1000$ قطب، ولتدفق النبضات؟ وأي طاقة للبت كان التدفق الخام سيتطلب؟
\end{exercise}

\begin{exercise}[\lvlA]
\label{ex:17:2}
\emph{كم بتًا في مئة عصبون.} يجمع المفكِّك نبضات $N$ عصبونًا في نوافذ $50\ms$ عند $r=20$ نبضة/ثانية، والارتباط الزوجي للضجيج $c=0.1$، والإشارة تغيّر تردد المجموعة بـ$30\%$ (جذر متوسط المربعات). قدّر المعدل $\tfrac12\log_2(1+\text{SNR})$ في النافذة عند $N=10$ و$100$ و$1000$ و$N\to\infty$. وعند $c=0$ و$N=1000$؟
\end{exercise}

\begin{exercise}[\lvlB]
\label{ex:17:3}
\emph{سرعة الواجهة-لوحة المفاتيح.} تختار الواجهة مرة ونصفًا في الثانية حرفًا من $N=32$: المقصود باحتمال $P$، والأخطاء متساوية الاحتمال. كم بتًا في الثانية تنقل عند $P=0.99$ و$0.94$ و$0.8$؟ وكم حرفًا في الثانية يلزم عند $P=0.94$ لبلوغ $10$ بتات/ثانية؟
\end{exercise}

\begin{exercise}[\lvlB]
\label{ex:17:4}
\emph{نمذجة: مفكِّك المجموعة.} $N$ عصبونًا بـ$\lambda_i=[b+m\,\mathbf v\cdot\mathbf p_i]_+$، $b=20$، $m=15$ نبضة/ثانية، نوافذ $50\ms$، و$\mathbf v$ بانحراف معياري $0.5$ لكل مركّبة؛ يرى الجميع ضجيجًا مشتركًا، $\mathbf v+\boldsymbol\xi$، $\sigma_\xi=0.2$. حاكِ مفكِّكًا غير متحيز بمتجه المجموعة. عند أي $N$ يتساوى الضجيجان، وكم بتًا للمركّبة في النافذة يعطي $N\to\infty$؟
\end{exercise}

\begin{exercise}[\lvlB]
\label{ex:17:5}
\emph{متعلمان.} يُخرج الدماغ الأمر $m=b\,v^*$، والمفكِّك $v=d\,m$، $\langle v^{*2}\rangle=1$. كلاهما بعد كل محاولة يخطو خطوة تدرج على $\tfrac12(v-v^*)^2$ بمعدل $\eta$. حاكِ بدءًا من $b=1.2$ و$d=1$ عند $\eta=0.8$ و$1.1$ و$1.5$ و$2$. كل منهما وحده مستقر عند $\eta<2$؛ فعند أي $\eta$ يستقر الزوج؟
\end{exercise}

\begin{exercise}[\lvlB]
\label{ex:17:6}
\emph{تصميم: خط معالجة المسبار.} $1024$ قناة بـ$20$ ألف عيّنة في الثانية، والعصبونات $10$ نبضات/ثانية، والإرسال اللاسلكي $1\,\text{nJ}$ للبت. أي عتبة على الضجيج الأبيض للعيّنات تعطي $0.1\,\text{Hz}$ على الأكثر من النبضات الكاذبة لكل قناة؟ نرسل أعداد النبضات في $20\ms$ بـ$2$ بت: ما القدرة، وكم مرة هي أصغر مما تتطلبه العيّنات الخام بـ$10$ بتات؟
\end{exercise}

\begin{exercise}[\lvlC]
\label{ex:17:7}
\emph{حد سرعة الواجهة.} قدّر $R\le DW\log_2(1+\text{SNR})$ لـ$D=10$ متغيرات بعرض نطاق $W=5\,\text{Hz}$ عند $\text{SNR}\approx0.9$ (ضجيج يشبه الإشارة) و$90$ (ضجيج مستقل لـ$1000$ عصبون). المُقاس نحو $10$ بتات/ثانية. أي تجربة تميّز بين تفسيري الفجوة: ضجيج يشبه الإشارة، أم الجهل بالشفرة؟
\end{exercise}
